\chapter{Theoretical Background}

In this chapter we present the mathematical and algorithmic foundations of the
LE-PSI system. All definitions follow standard notation from the lattice
cryptography literature. The chapter serves three purposes: (1) to establish the
formal primitives on which the DKLLMR23 protocol is built; (2) to formally define
the Laconic Encryption primitive, including the \textsf{LaconicEnc} function whose
inputs and correctness condition are central to understanding the PSI protocol; and
(3) to present a rigorous complexity analysis of auxiliary space --- the central
engineering problem solved in Phase II of this project.

% ─────────────────────────────────────────────────────────────────────────────
\section{Ring Learning With Errors (Ring-LWE)}

\begin{definition}[Ring-LWE~\cite{lyubashevsky2010ideal}]\label{def:rlwe}
Let $n$ be a power of two, $q$ a prime modulus, and
$R_q = \mathbb{Z}_q[x]/(x^n + 1)$
the cyclotomic ring. Let $\chi_s$ and $\chi_e$ be distributions over $R_q$
concentrated on polynomials with small-norm coefficients (discrete Gaussian with
standard deviation $\sigma = 3.2$). The \textbf{Ring-LWE} $(n, q, \chi_s, \chi_e)$
problem is: given arbitrarily many samples
$(a_i,\; b_i = a_i \cdot s + e_i \bmod q)$
where $a_i \xleftarrow{\$} R_q$, $s \leftarrow \chi_s$, $e_i \leftarrow \chi_e$,
distinguish these samples from uniformly random pairs $(a_i, u_i)$ with
$u_i \xleftarrow{\$} R_q$.
\end{definition}

The hardness of Ring-LWE is supported by a worst-case to average-case reduction
from hard lattice problems~\cite{lyubashevsky2010ideal}, and the closely related
Module-LWE assumption now underpins the NIST-standardised ML-KEM (FIPS 203) and
ML-DSA (FIPS 204) schemes, providing strong confidence in its post-quantum security.

\paragraph{Security parameters used in this project:}
\begin{table}[H]
\centering
\begin{tabular}{@{}lccl@{}}
\toprule
\textbf{Parameter} & \textbf{$D = 256$ mode} & \textbf{$D = 2048$ mode} & \textbf{Significance} \\
\midrule
Ring dimension $n$ & 256  & 2048 & Security vs.\ speed \\
Modulus $q$        & $\approx 2^{58}$ & $\approx 2^{58}$ & Error magnitude \\
Std deviation $\sigma$& 3.2  & 3.2  & Error distribution \\
Quantum security   & $\approx 64$-bit & $\approx 128$-bit & Via BKZ analysis~\cite{albrecht2015hardness} \\
Ciphertext size    & $\approx 2.7$\,KB & $\approx 10.8$\,KB & Network cost \\
\bottomrule
\end{tabular}
\end{table}

\begin{remark}\label{rem:security-tradeoff}
\textbf{Security level trade-off.} The 64-bit quantum security level at $D = 256$ is
below the NIST PQC standard of 128-bit quantum security. It represents a deliberate
engineering trade-off: $D = 2048$ provides 128-bit security but yields ciphertexts
approximately $4\times$ larger and computations approximately $4\times$ slower.
Chapter~6 presents correctness validation at $D = 2048$, and scaling to full 128-bit
benchmarks is identified as the primary direction for future work in Chapter~7.
\end{remark}

% ─────────────────────────────────────────────────────────────────────────────
\section{Number Theoretic Transform (NTT)}

\begin{definition}[NTT]\label{def:ntt}
The \textbf{Number Theoretic Transform} is the analogue of the Fast Fourier
Transform over $\mathbb{Z}_q$. For a polynomial
$f = \sum_{i=0}^{n-1} a_i x^i$, the NTT computes:
\[
\hat{f}(k) = \sum_{i=0}^{n-1} a_i\, \omega^{ik} \bmod q,
\qquad k = 0, \ldots, n-1,
\]
where $\omega$ is a primitive $n$-th root of unity in $\mathbb{Z}_q$.
\end{definition}

NTT enables polynomial multiplication in $R_q$ in $O(n \log n)$ time rather than
$O(n^2)$. In our profiling results (Section~6.6), NTT operations account for
approximately 25\% of total execution time and are identified as the primary
target for GPU acceleration in future work, since NTT is highly data-parallel.

% ─────────────────────────────────────────────────────────────────────────────
\section{Merkle Trees}

\begin{definition}[Merkle Tree~\cite{merkle1988}]\label{def:merkle}
A \textbf{Merkle tree} over leaves $\ell_0, \ldots, \ell_{m-1}$ is a complete
binary tree where each leaf node stores a cryptographic hash of its data element and each
internal node stores the hash of its two children's concatenation. The root node is a
short cryptographic digest committing to all $m$ leaves.
\end{definition}

\begin{definition}[Merkle Witness]\label{def:witness}
A \textbf{Merkle witness} (authentication path) for leaf $i$ is the sequence of
$\lceil \log_2 m \rceil$ sibling node hashes along the path from $\ell_i$ to the
root. Verification of the witness reconstructs the root hash in $O(\log m)$ time
and $O(\log m)$ space, confirming that $\ell_i$ is a genuine leaf of the committed tree.
\end{definition}

\begin{remark}\label{rem:merkle-storage}
\textbf{Storage optimisation.} A na\"{i}ve implementation allocates a separate length-$m$
array for each level of the tree, yielding $O(m \log m)$ total storage. We instead store
the entire tree in breadth-first order as a single flat array of $2m$ nodes, reducing
storage to $O(m)$. This is a straightforward implementation choice arising from the standard
binary heap layout; it is not presented as a novel contribution but is included here
as it directly affects the memory accounting in Section 3.6.
\end{remark}

% ─────────────────────────────────────────────────────────────────────────────
\section{Laconic Encryption --- Formal Definition}

This section provides the formal definition of the Laconic Encryption (LE) primitive,
including the \textsf{LaconicEnc} function. This definition was absent from the original
DKLLMR23 paper's exposition at the system level and is made explicit here to support
a clear understanding of the PSI protocol in Chapter~4.

\begin{definition}[Laconic Encryption Scheme]\label{def:laconic-enc}
A Laconic Encryption scheme over database universe $\{0,1\}^L$ consists of four algorithms:

\begin{itemize}
    \item \textsf{\textbf{LE.Setup}}$(1^\lambda, m) \to \mathsf{pp}$: Given security
          parameter $\lambda$ and database size $m$, outputs public parameters $\mathsf{pp}$
          including the ring dimension $n$, modulus $q$, and Merkle tree depth $\lceil \log_2 m \rceil$.
          
    \item \textsf{\textbf{LE.KeyGen}}$(\mathsf{pp}, D) \to (\mathsf{pk}, \mathsf{sk})$: Given a
          database $D = (d_0, \ldots, d_{m-1}) \in (\{0,1\}^L)^m$, constructs a Merkle tree
          over SHA-256 hashes of $D$, generates a Ring-LWE key pair $(pk_i, sk_i)$ for each
          leaf $i$, and outputs the aggregated public key $\mathsf{pk} = (\text{Merkle root } r, \{pk_i\})$
          and secret key $\mathsf{sk} = \{sk_i\}$.
          
    \item \textsf{\textbf{LaconicEnc}}$(\mathsf{pk}, j, \mathsf{msg}) \to \mathsf{ct}$: Given
          the public key $\mathsf{pk}$, a target position $j \in \{0, \ldots, m-1\}$, and a
          message $\mathsf{msg} \in \{0,1\}^L$, the client:
          \begin{enumerate}
              \item Retrieves the Merkle witness $\pi_j = (h_{j,1}, \dots, h_{j,\lceil \log_2 m \rceil})$ --- the sibling hashes along the path from leaf $j$ to the root.
              \item Encodes $\mathsf{msg}$ as a polynomial $m_j \in R_q$.
              \item Samples a Ring-LWE ciphertext: $\mathsf{ct} = (c_0, c_1)$ where
                    $c_0 = a \cdot r_j + e_0$, $c_1 = pk_j \cdot r_j + e_1 + \lfloor q/2 \rfloor \cdot m_j$,
                    with $r_j, e_0, e_1 \leftarrow \chi_e$ and $a \xleftarrow{\$} R_q$ tied to the
                    Merkle path $\pi_j$.
              \item Outputs $\mathsf{ct}$ together with $\pi_j$.
          \end{enumerate}
          
    \item \textsf{\textbf{LE.Dec}}$(\mathsf{sk}, i, \mathsf{ct}, \pi) \to \mathsf{msg}$ or $\perp$: Given
          secret key $sk_i$ for position $i$, ciphertext $\mathsf{ct}$, and witness $\pi$:
          \begin{enumerate}
              \item Verifies that $\pi$ is a valid Merkle witness for position $i$ under root $r$. If not, outputs $\perp$.
              \item Computes the decryption candidate: $\mathsf{msg}' = \textsf{Round}(c_1 - sk_i \cdot c_0)$.
              \item If $\mathsf{msg}'$ is a valid encoding of a database element, outputs $\mathsf{msg}'$; otherwise outputs $\perp$.
          \end{enumerate}
\end{itemize}
\end{definition}

\textbf{Correctness condition.} \textsf{LE.Dec}($\mathsf{sk}, i,$ \textsf{LaconicEnc}($\mathsf{pk}, j, \mathsf{msg}), \pi_j) = \mathsf{msg}$ if and only if $i = j$ (i.e., the server's leaf index matches the client's target position), which occurs precisely when $D[i] = \mathsf{msg}$ --- the server's $i$-th record matches the client's query element.

\begin{remark}[PSI connection]
The correctness condition directly implements private set intersection: the client encrypts each of its elements at the position it would occupy in the server's Merkle tree. The server attempts decryption at every leaf position. Successful decryption at position $i$ signals that the client's element matches $D[i]$, i.e., the element is in the intersection. This two-round structure is formalised in Chapter~4.
\end{remark}

% ─────────────────────────────────────────────────────────────────────────────
\section{Gadget Decomposition (GInvMNTT)}

\begin{definition}[Binary Gadget Decomposition]\label{def:gadget}
For a polynomial $f \in R_q$ with $q \approx 2^k$, the \textbf{binary gadget
decomposition} $\mathbf{G}^{-1}(f) \in R_2^k$ maps $f$ to a vector of $k$
binary polynomials satisfying:
\[
\sum_{j=0}^{k-1} 2^j \, [\mathbf{G}^{-1}(f)]_j = f \bmod q
\]
\end{definition}

\begin{remark}[Cost of gadget decomposition]\label{rem:gadget-cost}
At $q \approx 2^{58}$, each polynomial in $R_q$ expands into $k = 58$ binary polynomials
under $\mathbf{G}^{-1}$. Applied to a witness matrix of dimension $D \times D$, GInvMNTT
expands each witness by a factor of $58\times$. This expansion is cryptographically necessary:
it implements the LWE gadget required for correct decryption when Merkle path hashes are
incorporated into the ciphertext structure. It cannot be reduced without replacing the
underlying Ring-LWE construction. This $58\times$ factor is the root cause of the memory
scaling problem characterised in Section~3.6 and is also the principal source of the
performance gap between our Ring-LWE implementation and the pairing-based ALOS22 (discussed in Section~6.8).

Concretely: at $D = 256$, one witness pair occupies approximately 35~MB after gadget decomposition.
At $D = 2048$, one witness pair occupies approximately 280~MB.
\end{remark}

% ─────────────────────────────────────────────────────────────────────────────
\section{Formal Complexity Analysis}

\begin{definition}[Auxiliary Space Complexity]\label{def:auxspace}
For an algorithm $\mathcal{A}$ on input of size $N$, the
\textbf{auxiliary space complexity} $S_{\text{aux}}(\mathcal{A}, N)$ is the
maximum memory used by $\mathcal{A}$ \textit{excluding} the input and output.
\end{definition}

\subsection{Time Complexity}

\begin{table}[H]
\centering
\caption{Time complexity per phase of the LE-PSI protocol}
\label{tab:time-complexity}
\begin{tabular}{@{}llll@{}}
\toprule
\textbf{Phase} & \textbf{Dominant Operation} & \textbf{Time Complexity} & \textbf{Note} \\
\midrule
Key generation     & Ring-LWE sample  & $O(m \log m)$ & NTT, parallelisable \\
Tree construction  & Node hash     & $O(m \log m)$ & Breadth-first \\
Witness generation & GInvMNTT     & $O(m \cdot D^2 \log q)$ & Memory-bandwidth bound \\
Client encryption  & Ring-LWE encrypt & $O(n \log D)$ & CPU-bound \\
Intersection       & Decrypt check & $O(nm \cdot D)$ & Parallelised \\
\midrule
\textbf{Total} & & \textbf{$O(nmD)$} & \textbf{Intersection dominates} \\
\bottomrule
\end{tabular}
\end{table}

\subsection{Naïve Auxiliary Space}

\begin{theorem}[Naïve Auxiliary Space]\label{thm:naive-space}
The na\"{i}ve LE-PSI implementation, which generates all $m$ witness pairs and
holds them in RAM simultaneously before beginning intersection detection, has
auxiliary space complexity:
\[
S_{\text{aux}}^{\text{naive}}(m) = \Theta\!\left(m \cdot D^2 \cdot \log q\right).
\]
\end{theorem}

\begin{proof}
Each witness pair $(w^{(1)}_i, w^{(2)}_i)$ consists of two $D \times D$ matrices
of polynomials in $R_q$. After GInvMNTT, each polynomial expands into $\log_2 q$
binary polynomials of degree $D$. Per witness pair:
\[
\text{space} = 2 \times D^2 \times \log_2 q \times D \times 8\,\text{bytes}
             = \Theta(D^3 \log q)\,\text{bytes}.
\]
Storing all $m$ witnesses simultaneously:
$S_{\text{aux}}^{\text{naive}} = \Theta(m \cdot D^3 \cdot \log q)$.
Absorbing the constant $D$ factor into the asymptotic coefficient gives
$\Theta(m \cdot D^2 \cdot \log q)$ in ring elements.
\end{proof}

\begin{remark}\label{rem:naive-worst-case}
This worst case applies when all $m$ witnesses are generated and held in memory
simultaneously (fully-parallel execution). A single-threaded sequential implementation
reuses a single witness buffer of $\approx 35$~MB but requires $m$ times longer execution.
Our batched approach addresses both issues by bounding live witnesses to $B$ at any one
time, with $B \ll m$.
\end{remark}

\subsection{Batched Auxiliary Space}

\begin{theorem}[Batched Auxiliary Space]\label{thm:batch-space}
The batched LE-PSI implementation, which processes $B$ server records per
batch and discards witnesses before loading the next batch, has auxiliary
space complexity:
\[
S_{\text{aux}}^{\text{batch}}(B,m) =
  \underbrace{\Theta\!\left(B \cdot D^2 \cdot \log q\right)}_{\text{witness buffer, freed per batch}}
+ \underbrace{\Theta\!\left(m \cdot D^2\right)}_{\text{Merkle tree in RAM}}.
\]
\end{theorem}

\begin{proof}
During batch $b = \lfloor i/B \rfloor$, the algorithm allocates witnesses for
records $[bB,\, (b+1)B)$. These are freed before batch $b+1$ starts. Peak
witness allocation is $B \times \Theta(D^2 \log q)$. The in-memory Merkle tree holds $O(2m)$
polynomial nodes (without GInvMNTT expansion) at $\Theta(D)$ coefficients each,
totalling $\Theta(mD^2)$ coefficients, held throughout execution.
\end{proof}

\begin{corollary}\label{cor:batch-reduction}
Setting $B = O(1)$ makes the witness buffer $O(1)$ with respect to $m$.
Total auxiliary space is then $\Theta(m \cdot D^2)$, a reduction by
a factor of $\log_2 q \approx 58$ compared to the na\"{i}ve approach.
\end{corollary}

\subsection{Summary}

\begin{table}[H]
\centering
\caption{Auxiliary space complexity: naïve vs.\ batched}
\label{tab:complexity-summary}
\begin{tabular}{@{}lcc@{}}
\toprule
\textbf{Metric} & \textbf{Naïve} & \textbf{Batched (this work)} \\
\midrule
Time              & $O(nmD)$                       & \textbf{$O(nmD)$} --- unchanged \\
Aux.\ space (witnesses) & $\Theta(m \cdot D^2 \log q)$   & \textbf{$\Theta(B \cdot D^2 \log q)$} \\
Aux.\ space (tree)      & $\Theta(m \cdot D^2)$          & \textbf{$\Theta(m \cdot D^2)$} \\
\textbf{Total aux.\ space}& $\Theta(m \cdot D^2 \log q)$   & \textbf{$\Theta(m \cdot D^2)$} \\
Memory reduction  & ---                            & \textbf{Factor of $\log_2 q \approx 58\times$} \\
\bottomrule
\end{tabular}
\end{table}

The key insight is that batching trades parallelism for memory, while the in-memory
Merkle tree (Section~4.3) restores the parallelism by eliminating I/O latency during
witness generation. The combined effect --- batching plus parallelism --- is characterised
empirically in the ablation study of Section~6.7.

\clearpage
